\documentclass[14pt, a4paper]{extarticle}

\usepackage{indentfirst}
\usepackage{subfigure}
\usepackage{float}
\usepackage{multirow}
\usepackage{graphicx}
\usepackage[utf8]{inputenc}
\usepackage[T2A]{fontenc}
\usepackage[russian]{babel}
\usepackage{tempora}
\usepackage{ulem}
\usepackage{amsmath}
\usepackage{setspace}
\usepackage{amssymb}
\usepackage{tikz}
\usepackage{url}
\usepackage{titlesec}
\usepackage{array}
\usepackage{soul}

\usepackage[
    a4paper,
    top=25mm,
    bottom=25mm,
    left=20mm,
    right=20mm
]{geometry}

\setlength{\parindent}{1.25cm}
\linespread{1.2}
\sloppy
\pagestyle{empty}

\begin{document}

\newgeometry{left=20mm, right=20mm, top=20mm, bottom=20mm}
\thispagestyle{empty}
\selectlanguage{russian}

\begin{center}
    \begin{spacing}{1}
        \large\textbf{МОСКОВСКИЙ АВТОМОБИЛЬНО-ДОРОЖНЫЙ} \\
        \large\textbf{ГОСУДАРСТВЕННЫЙ ТЕХНИЧЕСКИЙ УНИВЕРСИТЕТ} \\ [0.1cm]
        \textbf{(МАДИ)}
    \end{spacing}
\end{center}

\vspace{-1.5cm}

\noindent
\begin{minipage}[t]{0.70\textwidth}
    \vspace{0pt}
    \raggedright
    Кафедра \\
    \underline{ИМПСИИ} \\ [0.2cm]
    Шифр и направление подготовки/специальность \\
    \underline{01.03.04 Прикладная математика} \\ [0.2cm]
    Направленность (профиль) \uline{Прикладная математика}

\end{minipage}%
\begin{minipage}[t]{0.30\textwidth}
    \vspace{5cm}
    \begin{flushright}
        \begin{minipage}{\textwidth}
            \raggedleft
            УТВЕРЖДАЮ \\
            зав. кафедрой \\ [0.6cm]
            \underline{\hspace{1.8cm}} / Яшина М.В. / \\ [0.4cm]
            «\underline{\hspace{0.7cm}}» \underline{\hspace{2.2cm}} 20\underline{\hspace{0.5cm}} г.
        \end{minipage}
    \end{flushright}
\end{minipage}

\vspace{1cm}
\begin{center}
    ЗАДАНИЕ \\
    НА ВЫПУСКНУЮ КВАЛИФИКАЦИОННУЮ РАБОТУ \\ [0.3cm]
    Чарыков Данила Владимирович \\
    \small{(Фамилия, Имя, Отчество обучающегося)}
\end{center}

\vspace{0.4cm}

\noindent Тема ВКР: \uline{Криптографические свойства булевых функций и их преобразований} \\
утверждена приказом МАДИ от «\underline{\hspace{0.6cm}}» \underline{\hspace{3cm}} 20\underline{\hspace{0.6cm}} г. № \underline{\hspace{2cm}}


\vspace{0.3cm}

\noindent 1. Исходные данные по ВКР: \\
\uline{Базовая организация -- Московский автомобильно-дорожный Государственный %
Технический Университет (МАДИ), Прикладная математика. Характер работы -- %
научно-исследовательская работа.}

\vspace{0.3cm}

\noindent 2. Обоснование темы ВКР и перечень подлежащих разработке вопросов: \\
\noindent \uline{В связи с возрастающими требованиями к информационной безопасности
и необходимостью защиты данных от современных киберугроз одним из актуальных вопросов
является автоматизация анализа криптографических свойств булевых функций, определяющих
стойкость алгоритмов шифрования.}


\newpage

\begin{table}[H]
\centering
\small
\begin{spacing}{1.1}
\begin{tabular}{|>{\centering\arraybackslash}p{1cm}|
                 >{\raggedright\arraybackslash}p{5cm}|
                 >{\centering\arraybackslash}p{4cm}|
                 >{\centering\arraybackslash}p{2cm}|
                 >{\centering\arraybackslash}p{1.5cm}|
                 >{\centering\arraybackslash}p{1.5cm}|}

\hline
\multirow{3}{*}{\centering № п/п} & \centering Наименование этапа работы & \centering Ф.И.О., должность, уч. степень, звание руководителя/консультанта & \centering Срок выполнения этапа, & \multicolumn{2}{c|}{подпись, дата} \\ \cline{5-6} 
 & \centering (раздела) & \centering (этапа работы, раздела) & \centering раздела & Задание выдал & Задание принял \\ \hline
1 & \centering Анализ и систематизация теоретических основ булевых функций & \centering Жуков Д.А. \\ к.ф.-м.н., доц. & 01.02.2025 & & \\ \hline
2 & \centering Разработка алгоритмического обеспечения и программная реализация приложения & \centering Жуков Д.А. \\ к.ф.-м.н., доц. & 01.04.2025 & & \\ \hline
3 & \centering Написание экономического обоснования & \centering Самохвалова Ж.П. \\ ст.пр. & 14.04.2025 & & \\ \hline
4 & \centering Написание текста диплома & \centering Жуков Д.А. \\ к.ф.-м.н., доц. & 01.05.2025 & & \\ \hline
5 & \centering Оформление текста ВКР & \centering Жуков Д.А. \\ к.ф.-м.н., доц. & 01.06.2025 & & \\ \hline
\end{tabular}
\end{spacing}
\end{table}

\vspace{1.5cm}

\noindent
\begin{tabular}{@{}p{3.5cm} >{\centering\arraybackslash}p{5.0cm} c >{\centering\arraybackslash}p{5.0cm}@{}}
Руководитель ВКР & \underline{\hspace{0.8cm}Жуков Д.А.\hspace{0.8cm}} & / & \underline{\hspace{5.0cm}} \\
& \scriptsize{(Ф.И.О.)} & & \scriptsize{(подпись, дата)} \\ [1cm]
Задание принял & & & \\
к исполнению & \underline{\hspace{0.8cm}Чарыков Д.В.\hspace{0.8cm}} & / & \underline{\hspace{5.0cm}} \\
& \scriptsize{(Ф.И.О.)} & & \scriptsize{(подпись, дата)}
\end{tabular}

\restoregeometry

\end{document}